% --- Άσκηση 38792 (38792.tex) ---
\Askhsh{38792}{4}
\begin{ylh}{2}
    \item 5.3. Γεωμετρική πρόοδος
\end{ylh}
Ένα είδος γρήγορα αναπτυσσόμενων τοξικών φυκιών εμφανίζεται την 1η Απριλίου, σε μια λίμνη ως αποτέλεσμα της προβληματικής διαχείρισης της γύρω περιοχής (κοντινά αγροκτήματα με απόβλητα και ένα εγκαταλελειμμένο ορυχείο υδραργύρου). Τα φύκια εμφανίζονται αρχικά σε μία μικρή ποσότητα $a$, χωρίς να γίνουν αμέσως αντιληπτά. Αρχίζουν να αναπτύσσονται και να καλύπτουν την επιφάνεια της λίμνης με τέτοιο τρόπο ώστε κάθε μέρα να καλύπτουν τη διπλάσια επιφάνεια από εκείνη που κάλυπταν την προηγούμενη ημέρα. Εάν συνεχίσουν να αναπτύσσονται με αυτόν τον ρυθμό, η επιφάνεια $\Lambda$ της λίμνης θα καλυφθεί πλήρως στις 30 Απριλίου, θέτοντας σε κίνδυνο όποιον οργανισμό ζει κοντά ή μέσα στη λίμνη.

\begin{alist}
\item Να γράψετε μία σχέση που να περιγράφει την επιφάνεια $\Lambda_{9}$ της λίμνης που έχει καλυφθεί από τα φύκια στις 9 Απριλίου και μία σχέση που να περιγράφει την επιφάνεια της λίμνης που έχει καλυφθεί από τα φύκια οποιαδήποτε ημέρα $\nu$ του Απριλίου.
\monades{6}
\item Ποιο μέρος της λίμνης έχει καλυφθεί από τα φύκια στις 26 Απριλίου; Να εξηγήσετε πώς καταλήξατε στην απάντησή σας.
\monades{6}
\item Ποια ημέρα τα φύκια θα έχουν καλύψει τη μισή λίμνη; Να εξηγήσετε πώς καταλήξατε στην απάντησή σας.
\monades{4}
\item Στις 29 Απριλίου βιολόγοι φτάνουν στη λίμνη και προσπαθούν να την καθαρίσουν από τα φύκια, όμως δεν καταφέρνουν να τα αφαιρέσουν εντελώς όλα. Εκτιμούν ότι το 0$,5\
